%%% FIELD proof-prop-conditional-polynomial-saturation
Put \(L=\operatorname{len}(\mathcal O)\), the encoding length of the immutable oracle input.  In one regime there are at most \(\binom D2\) observed--observed adjacencies, and each such adjacency has two reported endpoints.  Since the regimes are indexed by \(k=0,\ldots,K\), the number \(N_{\mathrm{end}}\) of reported endpoint positions therefore satisfies
\[
N_{\mathrm{end}}
 \le 2(K+1)\binom D2
 = (K+1)D(D-1)
 \le 2(K+1)D(D-1).
\tag{1}
\]
This count uses \(U_0=V\) and, for \(k\in[K]\), reports only the observed--observed endpoints in \(U_k=V\cup\{\zeta_k\}\), as prescribed by the endpoint domain in \(\mathrm{def:canonical\mbox{-}si\mbox{-}pag}\).

By the hypothesis on conclusions, every strict rule application replaces a circle at one reported endpoint by either \(\mathtt{tail}\) or \(\mathtt{arrow}\).  The same hypothesis says that this endpoint can never change strictly again.  Charging a strict application to an endpoint that it changes consequently gives
\[
N_{\mathrm{upd}}\le N_{\mathrm{end}}le 2(K+1)D(D-1).
\tag{2}
\]
Thus the asserted update bound follows directly from \((1)\).

It remains to account for the work needed to find those updates.  The candidate \(\mathcal R_{\mathrm{SI}}\) is fixed and finite, and every schema in the grammar \(\Lambda_{\mathrm{SI}}\) has finitely many typed vertex and regime placeholders by \(\mathrm{def:rule\mbox{-}language}\).  Hence there are constants \(c,q,r\in\mathbb N_0\), depending only on the fixed rule list and not on \(D\), \(K\), or \(\mathcal O\), for which the total number of instantiated rules in one exhaustive scan is at most
\[
N_{\mathrm{inst}}\le cD^q(K+1)^r.
\tag{3}
\]
If the rule list is empty, one may take \(c=1\) and \(q=r=0\).  By the premise-decider hypothesis and finiteness of the fixed list, there is a single polynomial \(p\) such that deciding the premise and every path predicate of any instance costs at most \(p(D,K,L)\).  Therefore an exhaustive scan costs at most
\[
cD^q(K+1)^r p(D,K,L),
\tag{4}
\]
which is polynomial in \(D\), \(K\), and \(L\).

Consider the deterministic saturation procedure that scans the instances in any fixed order, applies the first available strict conclusion, restarts the scan, and halts when a complete scan produces no change.  It performs at most \(N_{\mathrm{upd}}+1\) complete scans.  Combining \((2)\) and \((4)\), its work is bounded above by
\[
\bigl(2(K+1)D(D-1)+1\bigr)cD^q(K+1)^r p(D,K,L),
\tag{5}
\]
a polynomial in the stated input parameters.

The initialization in this procedure is exactly the initialization specified in \(\mathrm{def:monotone\mbox{-}closure}\).  Moreover, \(\mathrm{def:rule\mbox{-}language}\) makes the rules monotone: their conclusions only refine reported circles.  At termination no rule instance has an available strict conclusion, so the terminal state is closed.  Conversely, if \(Q\) is any closed state extending the same initialization, induction over the performed updates shows that every mark produced by the procedure is also present in \(Q\): the induction hypothesis embeds the current state in \(Q\), monotonicity preserves the premise of the selected rule instance in \(Q\), and closedness of \(Q\) supplies its conclusion.  The terminal state is therefore contained in every closed extension of the initialization.  It is the least monotone fixed point, hence it equals \(\operatorname{Cl}_{\mathcal R}(\mathcal O)\) by \(\mathrm{def:monotone\mbox{-}closure}\).

Every rule conclusion changes only a reported observed endpoint.  Thus, throughout the procedure, the observational skeleton and unshielded colliders on \(U_0=V\), and the complete skeletons and unshielded-collider families on \(U_k\) supplied by the faithful pooled selected laws for \(k\in[K]\), including facts involving \(\zeta_k\), remain immutable, again by \(\mathrm{def:rule\mbox{-}language}\) and \(\mathrm{def:monotone\mbox{-}closure}\).

Finally, suppose additionally that
\[
\operatorname{Cl}_{\mathcal R}(\mathcal O)
=\mathcal P_{\mathrm{SI}}(\mathcal O).
\tag{6}
\]
The left-hand side of \((6)\) is the polynomial-time output just computed, while the right-hand side is the canonical endpoint-invariance SI-PAG defined in \(\mathrm{def:canonical\mbox{-}si\mbox{-}pag}\).  Substitution in \((6)\) therefore identifies the output with that causal-structural target.  The procedure accesses only \(\mathcal O\), current endpoint marks, and rule predicates; at no step does it inspect or enumerate a member of \(\mathcal C(\mathcal O)\).  This proves every assertion of the proposition. \(\square\)
